\documentclass[11pt,a4paper]{report}

% ===================== PACKAGES =====================
\usepackage[utf8]{inputenc}
\usepackage[T1]{fontenc}
\usepackage[french]{babel}
\usepackage{amsmath,amssymb,amsthm}
\usepackage{graphicx}
\usepackage{geometry}
\geometry{margin=2.5cm}
\usepackage{hyperref}
\hypersetup{colorlinks=true, linkcolor=black, citecolor=black, urlcolor=blue}
\usepackage{booktabs}       % jolis tableaux
\usepackage{caption}
\usepackage{listings}       % pour les extraits de code
\usepackage{xcolor}
\usepackage{fancyhdr}
\usepackage{titlesec}
\usepackage[numbers,sort&compress]{natbib}

% ===================== MISE EN PAGE =====================
\pagestyle{fancy}
\fancyhf{}
\fancyhead[L]{\leftmark}
\fancyfoot[C]{\thepage}

% Style pour le code Python (lignes principales uniquement, cf. consignes)
\lstset{
    language=Python,
    basicstyle=\ttfamily\small,
    keywordstyle=\color{blue},
    commentstyle=\color{gray},
    stringstyle=\color{red},
    numbers=left,
    numberstyle=\tiny,
    frame=single,
    breaklines=true
}

% ===================== INFOS DOCUMENT =====================
\title{
    \Large Université de Strasbourg \\
    \large Master Mathématiques et Applications -- Parcours Statistique \\
    \vspace{1cm}
    \Huge \textbf{Titre du rapport de stage} \\
    \vspace{0.5cm}
    \large Sous-titre si besoin
}
\author{Yester MRTCHYAN}
\date{Année universitaire 2025--2026}

% ===================== DEBUT DOCUMENT =====================
\begin{document}

% ---------- Page de garde ----------
\maketitle
\thispagestyle{empty}

% ---------- Remerciements ----------
\chapter*{Remerciements}
\thispagestyle{empty}
% TODO : rédiger vos remerciements

\newpage

% ---------- Table des matières ----------
\tableofcontents
\thispagestyle{empty}

\newpage

% ===================== INTRODUCTION =====================
\chapter*{Introduction}
\addcontentsline{toc}{chapter}{Introduction}
% TODO : contexte du laboratoire, problématique, but du stage, annonce du plan

% ===================== PARTIE 1 =====================
\chapter{Contexte et cadre du stage}
L'Institut pluridisciplinaire Hubert Curien (IPHC) est un laboratoire de recherche du CNRS et de l'Université de Strasbourg. 
Differents axes de recherche existent dont le Département d'écologie, de physiologie et d'éthologie (DEPE). 
Ce dernier étudie les mécanismes permettant aux espèces animales de faire face aux contraintes de leur environnement à travers des approches d'écologie comportementale et d'écophysiologie, qu'elles soient fonctionnelles ou évolutives.
\section{Présentation du sujet de recherche}


\section{Planification du stage}
% TODO : chronologie annotation -> réentraînement -> matrices

\section{Principes du deep learning et du modèle d'identification utilisé}

\subsection{Apprentissage supervisé et réseaux de neurones}
% TODO : principe général, notion de fonction $\hat{f}_\theta$, fonction de perte

\subsection{Le modèle de classification utilisé}
% TODO : formaliser le problème

Le modèle prend en entrée une image $X_i$ et produit un vecteur de probabilités
\begin{equation}
    \hat{p}(X_i) = (\hat{p}_1, \hat{p}_2, \dots, \hat{p}_K), \qquad \sum_{k=1}^{K} \hat{p}_k = 1,
\end{equation}
où $K$ est le nombre d'individus connus dans la base. L'identité prédite est
\begin{equation}
    \hat{Y}_i = \arg\max_{k \in \{1,\dots,K\}} \hat{p}_k.
\end{equation}

\subsection{Ré-entraînement (fine-tuning)}
% TODO : pourquoi repartir du modèle existant plutôt que de zéro


% ===================== PARTIE 2 =====================
\chapter{Données et méthodologie d'annotation}

\section{Description de la base de données}
% TODO : structure des JSON, champs, volumétrie

\section{Protocole d'annotation et de correction manuelle}
% TODO

\section{Statistiques descriptives}

\subsection{Taux d'identification correcte}
% TODO : taux top-1, top-5, avant/après ré-entraînement

\subsection{Analyse des scores de confiance}
% TODO : distribution des scores selon exactitude de la prédiction

\section{Ré-entraînement du modèle : évaluation du gain}
% TODO


% ===================== PARTIE 3 =====================
\chapter{Construction des matrices de proximité sociale}

\section{Matrice estimée à partir des données annotées}
% TODO : définition mathématique

Soit $\hat{M}$ la matrice de proximité estimée, de taille $K \times K$, définie par
\begin{equation}
    \hat{M}_{ij} = \frac{\text{nombre de co-occurrences observées entre } i \text{ et } j}{\text{normalisation}}.
\end{equation}

\section{Matrice de référence}
% TODO : origine des données réelles

\section{Propriétés mathématiques des deux matrices}
% TODO : symétrie, positivité, diagonale, normalisation


% ===================== PARTIE 4 =====================
\chapter{Comparaison statistique des deux matrices}

\section{Problématique statistique}
% TODO : non-indépendance des entrées d'une matrice

\section{Présentation théorique de la méthode de comparaison}
% TODO : ex. test de Mantel -- statistique de test, hypothèse nulle, permutations

\section{Vérification des conditions d'application}
% TODO

\section{Résultats et interprétation}
% TODO


% ===================== PARTIE 5 =====================
\chapter{Discussion et limites}

\section{Synthèse croisée des résultats}
% TODO

\section{Limites de validité}
% TODO

\section{Apport personnel et autonomie}
% TODO


% ===================== CONCLUSION =====================
\chapter*{Conclusion}
\addcontentsline{toc}{chapter}{Conclusion}
% TODO


% ===================== BIBLIOGRAPHIE =====================
\addcontentsline{toc}{chapter}{Bibliographie}
\bibliographystyle{plainnat}
\bibliography{biblio}


% ===================== ANNEXES =====================
\appendix
\chapter{Annexes}

\section{Extraits de code}
% Exemple d'insertion de code (lignes principales uniquement, cf. consignes -- jamais le programme complet)
\begin{lstlisting}
# Exemple : calcul du taux d'identification correcte
correct = sum(1 for a in annotations if a["corrected_id"] == a["predicted_id"])
taux = correct / len(annotations)
\end{lstlisting}

\section{Captures d'écran de l'interface}
% TODO

\section{Tableaux complémentaires}
% TODO

\end{document}